% Part: applied-modal-logics
% Chapter: temporal-logic
% Section: introduction

\documentclass[../../../include/open-logic-section]{subfiles}

\begin{document}

\olfileid{aml}{tl}{int}

\olsection{పరిచయం}

కాలిక తర్కాలు భవిష్యత్తులో జరగబోయే లేదా గతంలో జరిగిన
విషయాల గురించిన వాదనలను పరిశీలిస్తాయి. కాలిక తర్కానికి
మూలకర్తగా ఆర్థర్ ప్రైయర్‌ను గుర్తిస్తారు; ఆయన దాన్ని
\emph{కాలరూప తర్కం} అని పిలిచారు. ఇక్కడి పరిశీలనలో
మనం ప్రధానంగా ప్రైయర్ మొదటి మోడల్ పద్ధతిని అనుసరిస్తాం:
సాధారణంగా కాలరూపం లేనివిగా వాదనలను చూసే
ప్రతిజ్ఞావాక్య తర్కపు ప్రాథమిక చట్రంలో కాలిక కారకాలను
ప్రవేశపెడతాం.

ఉదాహరణకు, ప్రతిజ్ఞావాక్య తర్కంలో బీజీ అనే కుక్క గురించి
మాట్లాడవచ్చు: కుక్కలకు అలవాటైనట్లే అది కొన్నిసార్లు
కూర్చుంటుంది, కొన్నిసార్లు కూర్చోదు. బీజీ కూర్చుంది,
అలాగే కూర్చోలేదు అని ఒకేసారి చెప్పడం సంప్రదాయ తర్కంలో
వైరుధ్యం. అయితే రెండు వాదనలూ సత్యం కావచ్చు---ఒకే
సమయంలో మాత్రం కాదు. భాషలో కాలిక కారకాలను చేర్చితే
దాన్ని సులభంగా వ్యక్తం చేయవచ్చు. ``బీజీకి బహుమతి
లేదా బంతి దొరుకుతుంది'' అనే వాక్యం నుంచి ``బీజీకి
బహుమతి దొరుకుతుంది లేదా బీజీకి బంతి దొరుకుతుంది''
అనే నిష్కర్ష చెల్లుబాటును కూడా అవి వివరించగలవు.

అయితే కాలిక తర్కంలో అనేక తాత్త్విక ప్రశ్నలు వస్తాయి;
వాటివల్ల ఒక కాలిక తర్క చట్రం బదులు మరొకదాన్ని
ఎంచుకోవచ్చు. ఉదాహరణకు, భవిష్యత్ అనిశ్చిత వాక్యం
భవిష్యత్తు గురించి చెబుతుంది; అది అనివార్యమూ కాదు,
అసాధ్యమూ కాదు. ``రిచర్డ్ రేపు కిరాణా దుకాణానికి
వెళ్తాడు'' అంటే, ఇంకా జరగని విషయం గురించి చెబుతున్నాం;
ఆ వాక్యపు సత్యమూల్యం వివాదాస్పదం. నిజానికి, దానికి
అసలు సత్యమూల్యం \emph{కేటాయించగలమా} అన్నదే
వివాదాస్పదం. మనం కఠిన నిర్ణయవాదులమైతే, ఆ సంఘటన
జరగాల్సిన సమయానికి ముందే ఆ వాక్యం సత్యమో అసత్యమో
అయి ఉంటుందని అంగీకరించవచ్చు---ఆ సత్యమూల్యం మనకు
ఇంకా తెలియకపోవచ్చు అంతే. దీనికి విరుద్ధంగా,
భవిష్యత్తు నిజంగా తెరచి ఉందనీ, భవిష్యత్ అనిశ్చిత
వాక్యాల సత్యమూల్యాలు ఇంకా నిర్ణయించబడలేదనీ నమ్మవచ్చు.

కాల నిర్మాణం, స్వభావం గురించిన ఈ అభిప్రాయాల్లో అనేకం
కాలిక తర్క నమూనాలు, చట్రాల ఎంపికలోనే ఇమిడి ఉంటాయి.
ఉదాహరణకు, కాలం రేఖీయంగా, శాఖలుగా విస్తరించేదిగా,
లేదా చక్రీయంగా ఉండే నమూనాల్లో ఏవి నిర్మించాలో
అడగవచ్చు. మన కాలిక నమూనాలకు ఆరంభ, ముగింపు బిందువులు
ఉండాలా, కాలాన్ని విడివిడి క్షణాలుగా లేదా నిరంతరంగా
చూపాలా అనే నిర్ణయాలూ తీసుకోవాలి.

\end{document}
